\subsection{Electronic supplementary tables}

The electronic result package contains the following thesis-facing tables:

\begin{itemize}
    \item the final gene-level evidence table for the seven comparative genes
    and OMD context;
    \item the 107-row candidate-level confidence and provenance table;
    \item the 103-row integrated protein/domain/MSA/SignalP table and per-gene
    conservation summaries;
    \item representative gene-structure and SynVoy gene-by-species evidence
    tables;
    \item mature-protein MEME/STREME/MAST summaries;
    \item promoter-sequence QC, corrected enrichment, and exploratory recurrence
    tables; and
    \item MGI/MP and HPO phenotype evidence summaries.
\end{itemize}

Accession, species, manual-decision, and source-provenance fields are retained so
that excluded or tentative candidates remain auditable rather than disappearing
from the final package.

\subsection{Electronic supplementary figures}

Supplementary figures include the seven gene-specific phylogenies and sequence
logos, a replicate-level GSE114919 expression display, eight representative
CDS-structure plots, ProtSpace projections, SynVoy
confidence/review summaries, the mature-protein motif-model matrix, promoter GC
and motif summaries, and the curated mouse phenotype heatmap. Individual
SignalP plots and tool-native motif reports are retained as diagnostics rather
than interpreted as independent biological results.

The original ZIP contained seven separate CDS-structure plots (BGN, EPYC, FMOD,
LUM, OGN, OMD, and PRELP). DCN was added later and is represented in the updated
eight-gene overview and full-transcript figures rather than by an obsolete or
newly invented per-gene panel.

\newcommand{\thesissuppfigure}[3]{%
\begin{figure}[p]
    \centering
    \includegraphics[width=\textwidth,height=0.80\textheight,keepaspectratio]{#1}
    \caption{#2}
    \label{#3}
\end{figure}%
}

\clearpage
\subsubsection{Expression context}

\thesissuppfigure
{figures/expression/expression_age_and_bgee_overview.png}
{Age-, zone-, species-, and Bgee-based expression context for the seven-gene
panel. The comparisons are descriptive and retain their dataset-specific
scales; missing Bgee calls are interpreted as missing database evidence rather
than absence of expression.}
{fig:supp-expression-context}

\thesissuppfigure
{figures/expression/gse114919_slrp_replicates.png}
{Replicate-level expression of the nine screened SLRPs in the deposited
GSE114919 mouse and rat growth-plate matrices. Points are biological replicates
and black bars are condition means on the authors' published normalized-value
scale. Conditions encode species, age, bone, and proliferative or hypertrophic
zone. Absolute values are not compared between species and no inferential test
is implied.}
{fig:supp-gse114919-replicates}

\clearpage
\subsubsection{Phylogenetic detail}

\thesissuppfigure
{figures/phylogeny/compact_per_gene_trees.png}
{Maximum-likelihood phylogenies of the seven curated gene-specific protein
sets. Models were selected independently and branch support was estimated with
1,000 SH-aLRT and 1,000 ultrafast-bootstrap replicates. All trees are unrooted;
their display orientation does not represent evolutionary direction.}
{fig:supp-per-gene-trees}

\thesissuppfigure
{figures/phylogeny/phylogenetic_result_overview.png}
{Phylogenetic result and provenance overview. The curated compact trees are the
source of thesis conclusions. Legacy very-large tree products are shown only as
an analysis audit because their former header parser could collapse unrelated
loci with repeated isoform labels; corrected locus-filtered inputs exist, but
their long-tree inference remains optional.}
{fig:supp-phylogeny-audit}

\clearpage
\subsubsection{Sequence-conservation logos}

\thesissuppfigure
{figures/protein/sequence_logos/BGN_sequence_logo.png}
{Occupancy-weighted amino-acid sequence logo for the curated BGN alignment.
Coordinates are alignment columns rather than human-protein residue numbers.}
{fig:supp-logo-bgn}

\thesissuppfigure
{figures/protein/sequence_logos/DCN_sequence_logo.png}
{Occupancy-weighted amino-acid sequence logo for the curated DCN alignment.
Coordinates are alignment columns rather than human-protein residue numbers.}
{fig:supp-logo-dcn}

\thesissuppfigure
{figures/protein/sequence_logos/EPYC_sequence_logo.png}
{Occupancy-weighted amino-acid sequence logo for the curated EPYC alignment.
Coordinates are alignment columns rather than human-protein residue numbers.}
{fig:supp-logo-epyc}

\thesissuppfigure
{figures/protein/sequence_logos/FMOD_sequence_logo.png}
{Occupancy-weighted amino-acid sequence logo for the curated FMOD alignment.
Coordinates are alignment columns rather than human-protein residue numbers.}
{fig:supp-logo-fmod}

\thesissuppfigure
{figures/protein/sequence_logos/LUM_sequence_logo.png}
{Occupancy-weighted amino-acid sequence logo for the curated LUM alignment.
Coordinates are alignment columns rather than human-protein residue numbers.}
{fig:supp-logo-lum}

\thesissuppfigure
{figures/protein/sequence_logos/OGN_sequence_logo.png}
{Occupancy-weighted amino-acid sequence logo for the curated OGN alignment.
The continuous conserved core is retained despite greater overall divergence.}
{fig:supp-logo-ogn}

\thesissuppfigure
{figures/protein/sequence_logos/PRELP_sequence_logo.png}
{Occupancy-weighted amino-acid sequence logo for the curated PRELP alignment.
Coordinates are alignment columns rather than human-protein residue numbers.}
{fig:supp-logo-prelp}

\clearpage
\subsubsection{Gene-structure detail}

\thesissuppfigure
{figures/gene_structure/full_gene_structure_5species.png}
{Full representative transcript organization for seven comparative genes and
OMD in human, mouse, cow, chicken, and zebrafish. Grey shows complete exon
extent, blue and orange denote derived UTR segments, and gene-specific colours
denote CDS. Coding structure is more stable than intron spacing and annotated
UTR length.}
{fig:supp-full-gene-structure}

\thesissuppfigure
{figures/gene_structure/coding_exon_splice_phase_map.png}
{Coding-exon boundaries and intron phases across five vertebrates. All 26
homologous coding junctions retain intron phase even where exact boundary
positions shift modestly.}
{fig:supp-splice-phase}

\thesissuppfigure
{figures/gene_structure/human_domain_exon_architecture.png}
{Human LRR-family Pfam hits projected onto coding exons. Domain-model overlap
with multiple exons links the protein architecture to the conserved coding
framework but does not imply that each exon is an independent functional
module.}
{fig:supp-domain-exon}

\thesissuppfigure
{figures/gene_structure/per_gene/BGN_cds_gene_structure_5species.png}
{Representative BGN CDS organization in five vertebrates. Boxes are CDS
segments and connecting lines are intervening genomic sequence in transcript
orientation.}
{fig:supp-structure-bgn}

\thesissuppfigure
{figures/gene_structure/per_gene/EPYC_cds_gene_structure_5species.png}
{Representative EPYC CDS organization in five vertebrates. Differences in
total span mainly reflect intervening sequence and transcript choice.}
{fig:supp-structure-epyc}

\thesissuppfigure
{figures/gene_structure/per_gene/FMOD_cds_gene_structure_5species.png}
{Representative FMOD CDS organization in five vertebrates.}
{fig:supp-structure-fmod}

\thesissuppfigure
{figures/gene_structure/per_gene/LUM_cds_gene_structure_5species.png}
{Representative LUM CDS organization in five vertebrates.}
{fig:supp-structure-lum}

\thesissuppfigure
{figures/gene_structure/per_gene/OGN_cds_gene_structure_5species.png}
{Representative OGN CDS organization in five vertebrates.}
{fig:supp-structure-ogn}

\thesissuppfigure
{figures/gene_structure/per_gene/OMD_cds_gene_structure_5species.png}
{Representative OMD CDS organization in five vertebrates. OMD is a context gene
and was not advanced through the complete seven-gene comparative workflow.}
{fig:supp-structure-omd}

\thesissuppfigure
{figures/gene_structure/per_gene/PRELP_cds_gene_structure_5species.png}
{Representative PRELP CDS organization in five vertebrates.}
{fig:supp-structure-prelp}

\clearpage
\subsubsection{Coding constraint and protein-space analyses}

\thesissuppfigure
{figures/evolutionary_rates/pairwise_dn_ds_human_reference.png}
{Human-reference pairwise coding-sequence constraint for seven comparative
genes and OMD. All 25 finite NG86 estimates are below one. Undefined or non-
positive synonymous distances are not plotted, and high-$d_S$ comparisons are
interpreted with saturation caution.}
{fig:supp-dnds}

\thesissuppfigure
{figures/protspace/protspace_canonical_projections.png}
{PCA, UMAP, and t-SNE projections of ProtT5 embeddings for the canonical
proteins and controls. These projections provide exploratory quality control;
visual proximity does not establish orthology or biological function.}
{fig:supp-protspace}

\clearpage
\subsubsection{Protein motifs}

\thesissuppfigure
{figures/motifs/protein_gene_motif_model_matrix.png}
{Compatibility of the 103 canonical proteins with the seven gene-specific
mature-protein MEME models. Every protein scores most strongly against its own
gene model, but the models were learned and evaluated on the same curated panel
and are not an independent accuracy estimate.}
{fig:supp-motif-matrix}

\thesissuppfigure
{figures/motifs/protein_candidate_motif_model_margins.png}
{Per-candidate margin between the assigned-gene motif model and the strongest
alternative model. Margins are supporting sequence-classification evidence and
must not be interpreted as validated functional binding motifs.}
{fig:supp-motif-margins}

\clearpage
\subsubsection{Promoter analysis}

\thesissuppfigure
{figures/regulatory/promoter_proximal_gc_heatmap.png}
{GC composition of strand-aware $-2{,}000/+200$-bp proximal promoters from 35
representative gene--species loci. All windows passed sequence-length and
ambiguity checks.}
{fig:supp-promoter-gc}

\thesissuppfigure
{figures/regulatory/promoter_targeted_tf_ame_enrichment.png}
{AME results for the targeted cartilage/osteogenesis transcription-factor
motif panel. HIF1A has the strongest raw result, but no targeted motif passes
the panel-corrected threshold; the figure does not demonstrate shared
regulation.}
{fig:supp-promoter-ame}

\thesissuppfigure
{figures/regulatory/promoter_targeted_tf_species_recurrence.png}
{Exploratory recurrence of targeted motif matches across five species at an
uncorrected FIMO threshold. The strict $q\leq0.05$ scan returned no sites, so
these patterns are hypotheses rather than conserved binding sites.}
{fig:supp-promoter-recurrence}

\clearpage
\subsubsection{Phenotype evidence}

\thesissuppfigure
{figures/phenotypes/mgi_single_gene_phenotype_heatmap.png}
{Breadth of curated single-gene mouse phenotype annotations for the seven
SLRPs. Counts represent database annotation coverage, not phenotype severity,
effect size, or cross-gene biological importance.}
{fig:supp-mgi-phenotypes}

\subsection{Manual-review records}

The manual sequence record documents the retained opossum and zebrafish OGN
proteins, tentative spotted-gar OGN and whale-shark EPYC proteins, and excluded
amphioxus SLRP-like candidate. The synteny review table uses the controlled
vocabulary \emph{accepted}, \emph{tentative}, \emph{rejected}, and
\emph{ambiguous}. All 98 rows have exact-assembly GFF/accession-based
decisions; the 25 prioritized exceptions received more detailed feature and
flanking-gene review. Only selected disputed cases were cross-checked visually
in NCBI Genome Data Viewer. The decisions integrate the target gene model,
informative flanking genes, reciprocal protein evidence, and tree placement.
